\textbf{Phần 1. Trắc nghiệm lựa chọn}
\setcounterpageref{ex}{0}

\begin{ex} %[2D8N42]
Trong phép tính $18 : 2 = 9$, số 18 được gọi là
\choice
{Số chia}
{\True Số bị chia}
{Thương}
{Số hạng}
\loigiai{
18 là Số bị chia.\\
\textbf{Đáp án của câu hỏi là B}
}
\end{ex}

\begin{ex} %[2D8N42]
Trong phép tính $18 : 2 = 9$, số 2 được gọi là
\choice
{\True Số chia}
{Số bị chia}
{Thương}
{Tích}
\loigiai{
2 là Số chia.\\
\textbf{Đáp án của câu hỏi là A}
}
\end{ex}

\begin{ex} %[2D8H42]
Trong phép tính $18 : 2 = 9$, số 9 được gọi là
\choice
{Tổng}
{Số chia}
{\True Thương}
{Số bị chia}
\loigiai{
Kết quả của phép chia gọi là Thương.\\
\textbf{Đáp án của câu hỏi là C}
}
\end{ex}

\begin{ex} %[2D8H42]
Tìm thương biết số bị chia là 30, số chia là 5:
\choice
{$5$}
{$7$}
{$8$}
{\True $6$}
\loigiai{
$30 : 5 = 6$.\\
\textbf{Đáp án của câu hỏi là D}
}
\end{ex}

\begin{ex} %[2D8H42]
Tìm thương biết số bị chia là 16, số chia là 2:
\choice
{$6$}
{\True $8$}
{$7$}
{$9$}
\loigiai{
$16 : 2 = 8$.\\
\textbf{Đáp án của câu hỏi là B}
}
\end{ex}

\begin{ex} %[2D8H42]
Số bị chia là 45, số chia là 5. Thương là
\choice
{\True $9$}
{$8$}
{$7$}
{$6$}
\loigiai{
$45 : 5 = 9$.\\
\textbf{Đáp án của câu hỏi là A}
}
\end{ex}

\begin{ex} %[2D8V42]
Thương của phép chia là 7, số chia là 2. Số bị chia là
\choice
{$12$}
{$16$}
{$9$}
{\True $14$}
\loigiai{
Số bị chia: $7 \times 2 = 14$.\\
\textbf{Đáp án của câu hỏi là D}
}
\end{ex}

\begin{ex} %[2D8V42]
Thương của phép chia là 8, số chia là 5. Số bị chia là
\choice
{\True $40$}
{$35$}
{$45$}
{$30$}
\loigiai{
Số bị chia: $8 \times 5 = 40$.\\
\textbf{Đáp án của câu hỏi là A}
}
\end{ex}

\begin{ex} %[2D8V42]
Số bị chia là 20, thương là 4. Số chia là
\choice
{$4$}
{$6$}
{\True $5$}
{$10$}
\loigiai{
Số chia: $20 : 4 = 5$.\\
\textbf{Đáp án của câu hỏi là C}
}
\end{ex}

\begin{ex} %[2D8V42]
Số bị chia là 18, thương là 9. Số chia là
\choice
{$1$}
{$3$}
{$4$}
{\True $2$}
\loigiai{
Số chia: $18 : 9 = 2$.\\
\textbf{Đáp án của câu hỏi là D}
}
\end{ex}

\begin{ex} %[2D8V42]
Tính giá trị biểu thức: Thương của 35 và 5 cộng với 18.
\choice
{$23$}
{$27$}
{\True $25$}
{$29$}
\loigiai{
$35 : 5 + 18 = 7 + 18 = 25$.\\
\textbf{Đáp án của câu hỏi là C}
}
\end{ex}

\begin{ex} %[2D8C42]
Tìm một số biết rằng lấy số đó chia cho 5 thì được kết quả là số chẵn liền sau của 6.
\choice
{$35$}
{\True $40$}
{$45$}
{$30$}
\loigiai{
Số chẵn liền sau của 6 là 8. Số cần tìm: $8 \times 5 = 40$.\\
\textbf{Đáp án của câu hỏi là B}
}
\end{ex}

\textbf{Phần 2. Trắc nghiệm đúng sai}
\setcounterpageref{ex}{0}

\begin{ex} %[2D8H42]
Xét tính đúng sai về thành phần phép chia:
\choiceTF
{\True Trong phép chia $25 : 5 = 5$, 25 là số bị chia.}
{\True Số 5 đứng sau dấu chia gọi là số chia.}
{Kết quả 5 gọi là tích.}
{\True Kết quả của phép chia gọi là thương.}
\loigiai{
\begin{itemize}
    \item 25 là số bị chia: Đúng.
    \item 5 là số chia: Đúng.
    \item Kết quả gọi là thương (không phải tích). Ý 3 Sai, Ý 4 Đúng.
\end{itemize}
\begin{itemize}
    \item a) \textbf{Đúng}.
    \item b) \textbf{Đúng}.
    \item c) \textbf{Sai}.
    \item d) \textbf{Đúng}.
\end{itemize}
}
\end{ex}

\begin{ex} %[2D8V42]
Cho phép chia có số bị chia là 16, số chia là 2. Xét tính đúng sai:
\choiceTF
{\True Thương của phép chia là 8.}
{\True Nếu gấp số bị chia lên 2 lần (thành 32) thì thương mới là 16.}
{Nếu giảm số bị chia đi 2 lần thì thương không đổi.}
{\True Khi đó thương mới là $8 : 2 = 4$.}
\loigiai{
\begin{itemize}
    \item Thương là 8: Đúng.
    \item $32 : 2 = 16$: Đúng.
    \item Thương giảm đi 2 lần. Ý 3 Sai, Ý 4 Đúng.
\end{itemize}
\begin{itemize}
    \item a) \textbf{Đúng}.
    \item b) \textbf{Đúng}.
    \item c) \textbf{Sai}.
    \item d) \textbf{Đúng}.
\end{itemize}
}
\end{ex}

\textbf{Phần 3. Trả lời ngắn}
\setcounterpageref{ex}{0}

\begin{ex} %[2D8H42]
Tìm thương biết số bị chia là 14, số chia là 2.
\shortans{7}
\loigiai{
\textbf{Đáp số:} $7$\\
$14 : 2 = 7$.
}
\end{ex}

\begin{ex} %[2D8H42]
Tìm thương biết số bị chia là 40, số chia là 5.
\shortans{8}
\loigiai{
\textbf{Đáp số:} $8$\\
$40 : 5 = 8$.
}
\end{ex}

\begin{ex} %[2D8V42]
Tìm số bị chia biết số chia là 2, thương là 9.
\shortans{18}
\loigiai{
\textbf{Đáp số:} $18$\\
$9 \times 2 = 18$.
}
\end{ex}

\begin{ex} %[2D8V42]
Tìm số bị chia biết số chia là 5, thương là 6.
\shortans{30}
\loigiai{
\textbf{Đáp số:} $30$\\
$6 \times 5 = 30$.
}
\end{ex}

\begin{ex} %[2D8V42]
Tìm số chia biết số bị chia là 50, thương là 10.
\shortans{5}
\loigiai{
\textbf{Đáp số:} $5$\\
$50 : 10 = 5$.
}
\end{ex}

\begin{ex} %[2D8C42]
Tìm thương của số tròn chục nhỏ nhất có hai chữ số chia cho 2.
\shortans{5}
\loigiai{
\textbf{Đáp số:} $5$\\
Số tròn chục nhỏ nhất là 10. Thương: $10 : 2 = 5$.
}
\end{ex}
